\documentclass[12pt]{article}
\usepackage[utf8]{inputenc}
\usepackage[T1]{fontenc}
\usepackage{hyperref}
\title{Tremembé Language Resources for K-12 Curriculum Development}
\author{Compiled by Claude Sonnet 5 (claude-sonnet-5) for Andr\'as Kornai}
\date{}
\begin{document}
\maketitle

\section{Introduction}
Tremembé (ISO 639-3 tme, Glottocode trem1235) is a dormant, unattested and unclassified language of the northern coast of Brazil, from Pará to Ceará. The Tremembé people live today mainly around Itarema and Itapipoca in Ceará. The searches found no revival programme for the language itself. They found language-policy and school-level vocabulary work. This is a thin report.

\section{Primary Educational Resources}
A school project document lists a subject called ``Pesquisa da Linguagem Tremembé'', alongside Torém and Tremembé history: \url{https://flacso.org.br/wp-content/blogs.dir/2184/files/2020/10/Caderno_de_Projetos_Indigena.pdf}

A 2025 article in REVASF studies language policy in Tremembé school education in Barra do Mundaú, Itapipoca, including vocabulary recovery: \url{https://periodicos.univasf.edu.br/index.php/revasf/article/download/3481/2170/16722}

\section{Linguistic and Reference Materials}
\begin{itemize}
\item Tremembé language (Wikipedia): \url{https://en.wikipedia.org/wiki/Trememb%C3%A9_language}
\item A 2015 dissertation from the Universidade Federal do Ceará, found in the search. I did not verify its content. \url{https://repositorio.ufc.br/ri/bitstream/riufc/20095/1/2015_dis_fpmachado.pdf}
\end{itemize}

\section{Religious and Community Resources}
Community identity centres on the Torém ritual dance rather than the language. A 2020 report covers the book \emph{Os Tremembé no Ceará: Tradição e Resistência}: \url{https://www.brasildefato.com.br/2020/10/29/livro-com-historias-de-lutas-do-povo-indigena-tremembe-e-lancado-no-ceara/}

\section{Digital and Audio Resources}
None were found.

\section{Curriculum Development Context}
There are no native speakers and almost no recorded language data. A curriculum would be a heritage and cultural-identity programme, led by the Tremembé community. It would use the school-based vocabulary recovery work above as its starting point, and it should not assume reconstructable grammar.
\end{document}